\documentclass{chsmc}
\chsmcsetup{
  year = 2019,
  part = 1,
  type = A,
  show-answers = false,
}
\begin{document}

\section{填空题：本大题共 8 小题，每小题 8 分，共 64 分.}

\begin{enumerate}
  \item 已知正实数 $a$ 满足 $a^a = (9a)^{8a}$，则 $\log_a 3a$ 的值为 \blankline
  \item 若实数集合 $\{1, 2, 3, x\}$ 的最大元素和最小元素之差等于该集合的所有元素之和，
    则 $x$ 的值为 \blankline
  \item 平面直角坐标系中，$\vec{e}$ 是单位向量，向量 $\vec{a}$ 满足
    $\vec{a}\cdot\vec{e} = 2$，且 $|\vec{a}|^2 \leqslant 5|\vec{a} + t\vec{e}|$
    对任意实数 $t$ 成立，则 $|\vec{a}|$ 的取值范围是 \blankline
  \item 设 $A$, $B$ 为椭圆 $\Gamma$ 的长轴顶点，$E$, $F$ 为 $\Gamma$ 的两个焦点，
    $|AB| = 4$，$|AF| = 2 + \sqrt{3}$，$P$ 为 $\Gamma$ 上一点，
    满足 $|PE|\cdot|PF| = 2$，则 $\triangle PEF$ 的面积为 \blankline
  \item 在 $1, 2, 3, \cdots, 10$ 中随机选出一个数 $a$，
    在 $-1, -2, -3, \cdots, -10$ 中随机选取一个数 $b$，
    则 $a^2 + b$ 被 $3$ 整除的概率是 \blankline
  \item 对任意区间 $I$，用 $M_I$ 表示函数 $y = \sin x$ 在 $I$ 上的最大值.
    若正数 $a$ 满足 $M_{[0,a]} = 2M_{[a,2a]}$，则 $a$ 的值为 \blankline
  \item 如图，在正方体 $ABCD\text{-}EFGH$ 的一个截面经过顶点 $A$, $C$ 及棱 $EF$
    上一点 $K$，且将正方体分成体积比为 $3:1$ 的两部分，
    则 $\dfrac{EK}{KF}$ 的值为 \blankline
  \item 将 $6$ 个数 $2, 0, 1, 9, 20, 19$ 按任意次序排成一行，
    拼成一个 $8$ 位数（首位不为 $0$），则产生的不同的 $8$ 位数的个数为 \blankline
\end{enumerate}

\section{解答题：本大题共 3 个小题，满分 56 分. 解答应写出文字说明、证明过程或演算步骤.}

\begin{enumerate}[resume]
  \item (本题满分 16 分) 在 $\triangle ABC$ 中，$BC = a$，$CA = b$，$AB = c$.
    若 $b$ 是 $a$ 与 $c$ 的等比中项，且 $\sin A$ 是 $\sin(B-A)$ 与 $\sin C$
    的等差中项，求 $\cos B$ 的值.
  \item (本题满分 20 分) 在平面直角坐标系 $xOy$ 中，圆 $\Omega$ 与抛物线
    $\Gamma\colon y^2 = 4x$ 恰有一个公共点，且圆 $\Omega$ 与 $x$ 轴相切于
    $\Gamma$ 的焦点 $F$. 求圆 $\Omega$ 的半径.
  \item (本题满分 20 分) 称一个复数列 $\{z_n\}$ 为“有趣的”，若 $|z_1| = 1$，
    且对任意正整数 $n$，均有 $4z_{n+1}^2 + 2z_n z_{n+1} + z_n^2 = 0$.
    求最大的常数 $C$，使得对一切有趣的数列 $\{z_n\}$ 及任意正整数 $m$，
    均有 $|z_1 + z_2 + \cdots + z_m| \geqslant C$.
\end{enumerate}

\end{document}